\chapter{总结与展望}

\section{总结}

本文完成了一个面向 PDF 与图片文档的多模态理解问答系统。系统从文档上传开始，依次完成页面解析、文本与视觉证据构建、向量检索、重排、GRPO 决策模型、答案生成和证据溯源展示，形成了从后端模型 pipeline 到前端 PDF 预览界面的完整闭环。用户上传文档后，系统能够将页面内容切分为 text、table、figure、code 等结构化 chunk，同时保留页码、bbox 和页面图路径，使问答结果可以回到原始 PDF 中进行高亮定位。

从系统实现看，本项目围绕文档问答构建了一套可检索、可定位、可验证的证据组织方式。文档解析模块负责把 PDF 或图片转化为可入库的文本和视觉区域；embedding 与检索重排模块负责召回相关 chunk 和页面区域；多模态生成模块根据问题类型组织局部视觉 crop、图表区域或文本上下文；前端通过连续 PDF 预览和区域高亮，把模型回答与原文档证据连接起来。因此，系统既能回答纯文本问题，也能处理图表数值、广告标签、票据版面和表格行列关系等需要视觉信息参与的问题。

实验结果也验证了多模态证据的价值。Our-RAG 的总体 EM 达到 0.9750，而 Text-RAG 的总体 EM 为 0.3750。这个差异说明，文档问答的提升并不只来自更长的文本上下文，而是来自更准确的证据组织：系统先通过检索和重排找到相关页面，再使用局部视觉区域补充 OCR 文本中缺失的版面、图表和视觉实体信息，最后由 VLM 结合证据生成答案。

从工程实现角度看，本文系统有三点主要设计：一是将文档解析、视觉区域、检索重排和 PDF 高亮统一到 PDF/图片文档问答任务中；二是使用统一的文本-图像 embedding 接口，使文本、表格和视觉证据可以进入同一检索流程；三是通过 GRPO 决策模型组织问题特征、候选 LLM/VLM 打分、证据组合和策略轨迹记录，使模型选择参与实际问答流程。整体来看，本文系统验证了多模态 RAG 在文档智能助手中的可行性。

\section{局限与未来工作}

当前系统已经形成从文档解析、检索重排、GRPO 决策模型、多模态生成到 PDF 证据高亮的完整闭环，但从系统模型能力和更长远的 AGI 视角来看，仍然存在一些值得继续推进的问题。

\begin{itemize}
  \item \textbf{系统模型局限}：当前系统仍然依赖外部 OCR、版面解析、向量检索、Reranker 和 VLM 等多个模块串联工作。每个模块的误差都会影响最终回答，例如 OCR 漏识别会导致 evidence 不完整，bbox 偏移会影响前端高亮，检索召回不足会使 VLM 看不到关键区域。因此，系统虽然具备多模态问答闭环，但本质上仍是模块化系统，而不是端到端统一理解文档的模型。
  \item \textbf{推理与记忆能力局限}：系统能够根据当前文档和当前问题完成检索问答，但对长程任务规划、跨文档持续记忆和多轮目标追踪的支持仍然有限。它可以回答“这个页面中的信息是什么”，但还不能像真正的智能助理一样长期维护用户目标、主动发现文档之间的矛盾，或者在多轮交互中形成稳定的任务计划。
  \item \textbf{AGI 视角下的潜力}：文档问答是通向更通用智能助手的重要场景之一。真实知识往往分布在 PDF、表格、图片、图表、扫描件和网页截图中，一个接近 AGI 的系统需要同时具备阅读、定位、比较、推理和溯源能力。本文系统已经把多模态感知、检索增强、模型选择和证据高亮结合在一起，说明文档智能可以作为连接外部知识与大模型推理能力的有效入口。
  \item \textbf{AGI 视角下的瓶颈}：当前系统的理解能力仍然受限于候选 evidence 的质量和底层模型的视觉推理能力。当文档存在复杂版面、跨页表格、隐含图表关系或需要常识推理时，系统仍可能出现证据定位正确但推理不充分、模型回答合理但来源不够严谨的问题。未来需要更强的多模态基础模型、更稳定的视觉检索、更可靠的自我验证机制，以及能够持续学习的决策模型。
\end{itemize}

整体来看，本文系统已经完成多模态文档问答所需的主要模块，并通过实验验证了视觉 evidence、Reranker、统一多模态向量化和 GRPO 决策模型的作用。后续工作可以从两个方向继续推进：一方面提升复杂文档中的解析、定位和生成稳定性；另一方面从 AGI 视角增强系统的长期记忆、任务规划、自我评估和持续学习能力，使其从“能回答文档问题的系统”进一步走向“能理解、使用并管理文档知识的智能体”。
